\section{結論}
\label{sec:conclusion}

本研究從容量理論角度分析 P2C 負載均衡下的慢故障緩解，
得出三項主要結論。
\emph{操作包絡線}（$\fmax = 1 - L/U$）決定緩解何時可行。
\emph{臨界嚴重度}（$\sstar = fq/((1-q)(N-f))$）
識別出哪些減速區間可由 P2C 的隱式迴避處理，哪些則需要顯式偵測。
\emph{SHED--P2C 互補性}（命題~\ref{thm:complementarity}）表明，在佇列感知路由下，
偵測速度比策略選擇更為關鍵：漸進式權重削減與完全隔離
達到幾乎相同的 P99，且任何顯式緩解的效益
隨叢集大小以 $O(N^{-1.3})$ 衰減。

透過涵蓋 15 個故障場景的評估，
本研究歸納出兩項具實務指導意義的發現：
（1）反應式機制（對沖請求、超時重試）對持續性慢故障無效；
（2）主動偵測與流量排除既必要又充分，
且在 P2C 下具體的排除機制並不重要。

嚴重度非單調性在實務上意義重大：
僅追蹤 P99 延遲的監控面板將無法偵測 $s > \sstar$ 的故障。
此類故障的 P99 延遲表現正常，
卻在 P999 造成災難性影響。
系統應監控多個百分位以偵測這些「沉默殺手」。

\subsubsection*{局限性}

本研究存在若干局限：模型假設為單層架構，
各工作節點同質且請求之間無親和性約束。
多層系統的扇出會放大慢故障影響~\cite{dean2013tail}，
預期操作包絡線將因此收緊。
利用率上限 $U$ 由平均延遲推導
（推論~\ref{cor:U}）；
對本文所使用的 M/D/1 模型而言，此公式為保守估計。
對於高變異服務時間（$\Cs^2 \geq 1$），
基於平均值的 $U$ 則為樂觀估計（備註~\ref{rem:mean-p99}）。
SHED--P2C 互補性的結論以 P2C 路由為前提。
在加權隨機等非佇列感知路由下，
SHED 仍能帶來顯著效益（表~\ref{tab:shed_complementarity}），故此結論不可直接推廣至其他路由架構。

\subsubsection*{未來工作}

未來可朝以下兩個方向延伸。
第一，\emph{重試放大}：慢故障下的重試會使有效負載增加為
$L_{\text{eff}} = L \cdot (1 + r \cdot p_{\text{timeout}})$，
可能縮減操作包絡線。
將重試與容量極限的交互作用形式化，
對具有重試策略的系統至關重要。
第二，將分析擴展至\emph{扇出架構}，
其中單一慢節點經由散集-聚集影響 $O(N)$ 個請求，
容錯能力可能降至
$\fmax \approx (1-L/U)/\text{fan-out}$。
